% SPDX-License-Identifier: PMPL-1.0-or-later
% SPDX-FileCopyrightText: 2026 Jonathan D.A. Jewell
%
% Graduated Trust: A Spectrum of Type Safety
% from Localhost to the Open Internet
%
% Intended venue: arXiv preprint (cs.NI / cs.CR / cs.PL)
% Subsequently: NSDI, USENIX Security, or IEEE S&P
%
% Third paper in the trilogy:
%   1. Groove Protocol (the mechanism)
%   2. Typing the Boundaries (the taxonomy)
%   3. Graduated Trust (the trust spectrum) ← this paper

\documentclass[11pt,a4paper]{article}
\usepackage[utf8]{inputenc}
\usepackage[T1]{fontenc}
\usepackage{amsmath,amssymb,amsthm}
\usepackage{listings}
\usepackage{hyperref}
\usepackage{booktabs}
\usepackage{graphicx}
\usepackage{xcolor}
\usepackage[margin=2.5cm]{geometry}

\definecolor{codegreen}{rgb}{0.25,0.49,0.31}
\definecolor{codegray}{rgb}{0.5,0.5,0.5}
\definecolor{codeblue}{rgb}{0.22,0.40,0.72}

\lstset{
  basicstyle=\ttfamily\small,
  keywordstyle=\color{codeblue}\bfseries,
  commentstyle=\color{codegray}\itshape,
  stringstyle=\color{codegreen},
  breaklines=true,
  frame=single,
  captionpos=b
}

\newtheorem{theorem}{Theorem}
\newtheorem{definition}{Definition}
\newtheorem{property}{Property}
\newtheorem{observation}{Observation}

\title{Graduated Trust: A Spectrum of Type Safety\\
from Localhost to the Open Internet}

\author{Jonathan D.A. Jewell\\
\texttt{j.d.a.jewell@open.ac.uk}\\
The Open University}

\date{March 2026}

\begin{document}

\maketitle

\begin{abstract}
Type safety in distributed systems is traditionally binary: either both
sides share a compiled schema (full safety) or they exchange untyped
bytes (no safety). We argue this dichotomy is false and present a
five-level graduated trust spectrum for typed communication, ranging
from compile-time proofs on localhost to runtime-validated declarations
on the open internet.

We introduce IPv6T, a typed payload framing protocol that embeds type
hashes, capability identifiers, and provenance chain links inside
standard TCP/IPv6 payloads. IPv6T requires no changes to network
infrastructure --- it passes through every router, NAT, firewall, CDN,
and proxy unmodified. Combined with the Groove Protocol for local
discovery and the Typed Frame Router for protocol translation, IPv6T
extends the type safety spectrum from localhost to the WAN.

We are honest about what IPv6T does and does not provide: on the
internet, types are \emph{attested}, not \emph{proven}. A type hash
declares intent; it does not constitute a compile-time proof. Runtime
validation is still required. The contribution is not ``types cross the
internet'' but rather: \emph{there exists a continuous spectrum from
untyped to proven, and we can push any communication channel upward on
that spectrum with known costs and guarantees at each level.}

We specify the 105-byte GRV6 frame format, define the five trust
levels, analyse the security properties and limitations of each level,
and demonstrate that testing the entire system is trivial --- five tests,
under 50 lines each, under 1 second total runtime.
\end{abstract}

\section{Introduction}

\subsection{The Binary Fallacy}

Current approaches to typed communication treat type safety as binary:

\begin{itemize}
\item \textbf{Full safety:} both sides compile against the same schema
  (Protocol Buffers, Cap'n Proto, FlatBuffers). The schema is shared at
  build time. Violations are compile errors.
\item \textbf{No safety:} data is exchanged as JSON, XML, or raw bytes.
  Types are reconstructed by parsing. Violations are runtime exceptions.
\end{itemize}

This binary framing forces a choice: tight coupling (shared schema) or
no safety (untyped exchange). Most real systems choose the latter because
tight coupling across the internet is impractical.

\subsection{The Spectrum}

We argue that type safety is not binary but graduated. Between ``proven
at compile time'' and ``completely untyped'' there are useful
intermediate levels with real benefits and honest limitations.

\begin{definition}[Trust Level]
A \emph{trust level} is a point on the spectrum from untyped to proven,
characterised by:
\begin{enumerate}
\item What the sender asserts about the data's type.
\item What evidence accompanies the assertion.
\item What the receiver must do to validate the assertion.
\item What guarantee holds if validation succeeds.
\end{enumerate}
\end{definition}

\section{The Five Trust Levels}

\begin{table}[h]
\centering
\small
\begin{tabular}{clllp{4cm}}
\toprule
Level & Name & Context & Evidence & Guarantee \\
\midrule
0 & Untyped & Raw internet & None & None. Parse and pray. \\
1 & Declared & GRV6 header & Type hash & Sender claims a type. Receiver validates at runtime. Early rejection on mismatch. \\
2 & Attested & GRV6 + prov & Type hash + provenance chain & Declared + auditable history. You can prove \emph{when} data was sent and \emph{what type was claimed}. \\
3 & Verified & Groove + GRV6 & Structural match + manifests & Both sides have Groove manifests. Types structurally matched at connection time. Not compiled together but validated at handshake. \\
4 & Proven & Idris2 + Groove & Compile-time proof & Both sides compiled against the same Idris2 ABI. The compiler verified compatibility. Zero runtime validation needed. \\
\bottomrule
\end{tabular}
\caption{The five trust levels}
\label{tab:levels}
\end{table}

\subsection{Level 0: Untyped}

The status quo. Data is sent as bytes over TCP. The receiver parses it
according to an implicit contract (a URL path, a Content-Type header, a
shared understanding between developers). If the contract is violated,
the error manifests at runtime --- often in production, often silently.

\textbf{Cost:} zero protocol overhead.\\
\textbf{Guarantee:} none.

\subsection{Level 1: Declared}

The sender prepends a 105-byte GRV6 frame header containing a SHA-256
hash of the A2ML type signature. The receiver compares the hash against
its expected type(s).

\textbf{What it buys:}
\begin{itemize}
\item Early rejection: mismatched traffic is rejected at the frame level
  before payload parsing begins. No CPU wasted on incompatible data.
\item Schema negotiation: two sides can compare type hashes before
  exchanging any application data. Fast ``do we speak the same type?''
  check.
\item Discoverability: a receiver seeing a GRV6 header knows the sender
  is type-aware, even without prior arrangement.
\end{itemize}

\textbf{What it does NOT buy:}
\begin{itemize}
\item No proof that the payload conforms to the declared type. The hash
  is an assertion, not evidence. Runtime validation is still required.
\item No proof that the sender is honest. Any actor can compute a type
  hash and attach it to arbitrary data.
\end{itemize}

\textbf{Cost:} 105 bytes per frame. Negligible for payloads $> 1$KB.\\
\textbf{Guarantee:} the receiver knows the sender \emph{intended} to
send data of this type.

\subsection{Level 2: Attested}

Level 1 plus a provenance chain. Each frame's \texttt{prov\_hash} links
to the previous frame, forming a hash chain. Optionally stored in
VeriSimDB for persistent audit.

\textbf{What it adds:}
\begin{itemize}
\item Auditability: you can verify the sequence, timing, and type
  declarations of all frames in a conversation.
\item Non-repudiation (weak): the sender cannot deny sending a frame
  without breaking the hash chain.
\item Compliance: regulatory contexts (financial, medical) that require
  data lineage get it for free.
\end{itemize}

\textbf{What it does NOT add:}
\begin{itemize}
\item Still no proof of payload conformance.
\item The provenance chain proves ordering and consistency, not
  correctness.
\end{itemize}

\textbf{Cost:} 105 bytes + provenance storage.\\
\textbf{Guarantee:} Level 1 + verifiable history.

\subsection{Level 3: Verified}

Both sides have Groove capability manifests. Before exchanging data,
they perform a Groove handshake: compare manifests, structurally match
offered and consumed capabilities, establish a connection.

\textbf{What it adds:}
\begin{itemize}
\item Structural compatibility: the receiver knows the sender's
  \emph{manifest} matches its requirements, not just that a hash
  matches. The full A2ML type definition was compared, not just a hash.
\item Bidirectional: both sides verified each other's capabilities.
\item Dynamic: if the sender's capabilities change (upgrade, downgrade),
  the next Groove handshake catches it.
\end{itemize}

\textbf{What it does NOT add:}
\begin{itemize}
\item Still runtime validation. The manifest comparison happened at
  connection time, but the payload could still be malformed (buggy
  sender, corrupted in transit, malicious modification after handshake).
\item The two sides were not compiled together. They share a type
  definition (via manifests) but not a compiler proof.
\end{itemize}

\textbf{Cost:} Groove handshake (one-time, \textasciitilde 5ms) + 105
bytes per frame.\\
\textbf{Guarantee:} Level 2 + structural capability match.

\subsection{Level 4: Proven}

Both sides were compiled against the same Idris2 ABI definitions. The
compiler has verified that the sender's output type and the receiver's
input type are structurally identical. The roundtrip proofs guarantee
encoding and decoding are lossless. The linear handle discipline
guarantees the connection is properly managed.

\textbf{What it adds:}
\begin{itemize}
\item Compile-time proof of compatibility. Not a runtime check --- a
  mathematical certainty.
\item The receiver can skip runtime validation (in trusted contexts)
  because the proof exists.
\item Resource safety: the connection handle is linear; leaks and
  double-closes are compile errors.
\end{itemize}

\textbf{When it applies:}
\begin{itemize}
\item Both sides compiled from the same source repository.
\item Both sides are on the same machine (localhost Groove).
\item Both sides are in the same build pipeline (CI/CD verification).
\end{itemize}

\textbf{When it does NOT apply:}
\begin{itemize}
\item The remote side is a third party. You did not compile their code.
\item The remote side is a modified binary. The proof may be voided.
\item The remote side is across the public internet and you have no
  build-time relationship.
\end{itemize}

\textbf{Cost:} zero runtime cost (proofs erased at compile time).\\
\textbf{Guarantee:} mathematical proof of type preservation.

\section{IPv6T: The Wire Protocol}

\subsection{Design Rationale}

We considered three approaches for carrying type information across the
internet:

\begin{enumerate}
\item \textbf{IPv6 extension headers:} technically correct but
  practically broken. Middleboxes drop packets with unknown extension
  headers at rates of 10--40\% \cite{google_eh2015}.
\item \textbf{IPv6 flow labels:} too small (20 bits) for meaningful type
  information.
\item \textbf{In-payload framing:} works through every middlebox on the
  internet because it is inside the TCP payload, which network
  infrastructure passes through unmodified.
\end{enumerate}

We chose approach 3. The GRV6 frame header is application data. The
internet carries it faithfully because the internet carries all
application data faithfully.

\subsection{Frame Format}

The GRV6 frame header is 105 bytes:

\begin{lstlisting}[caption={GRV6 frame header layout}]
Bytes  0-3:   Magic "GRV6" (0x47525636)
Byte   4:     Version (1)
Byte   5:     Flags [TYPED|PROVEN|ATTESTED|COMPRESSED]
Bytes  6-7:   Reserved (0x0000)
Bytes  8-39:  type_hash  (SHA-256 of A2ML type signature)
Bytes 40-71:  cap_hash   (SHA-256 of capability manifest)
Bytes 72-103: prov_hash  (SHA-256 of previous provenance)
Bytes 104-107: payload_length (uint32 big-endian)
Bytes 108+:   payload
\end{lstlisting}

\subsection{Why Not Just Use Protobuf?}

Protocol Buffers embed type information via field numbers. The schema is
shared at build time. This provides Level 4 (proven) safety when both
sides compile the same \texttt{.proto} file, and Level 0 (untyped)
otherwise. There is no intermediate.

GRV6 provides Levels 1--3: type awareness without requiring build-time
schema sharing. It is complementary to Protobuf, not a replacement.
A Protobuf payload can be wrapped in a GRV6 frame to add attestation
and capability hashing that Protobuf lacks.

\section{Security Analysis}

\subsection{What Each Level Defends Against}

\begin{table}[h]
\centering
\small
\begin{tabular}{lccccc}
\toprule
Attack & L0 & L1 & L2 & L3 & L4 \\
\midrule
Schema mismatch (accidental) & \texttimes & \checkmark & \checkmark & \checkmark & \checkmark \\
Malformed payload (buggy sender) & \texttimes & \texttimes & \texttimes & \texttimes & \checkmark \\
Type spoofing (malicious sender) & \texttimes & \texttimes & \texttimes & Partial & \checkmark \\
Replay attack & \texttimes & \texttimes & \checkmark* & \checkmark* & \checkmark \\
Provenance forgery & N/A & N/A & Hard & Hard & Impossible \\
Schema evolution breakage & \texttimes & \checkmark & \checkmark & \checkmark & \checkmark \\
\bottomrule
\end{tabular}
\caption{Defences by trust level. *Requires prov\_hash tracking.}
\end{table}

\subsection{Honest Limitations}

\begin{enumerate}
\item \textbf{Level 1 is an assertion, not a proof.} The type hash says
  ``I claim this is type X.'' It does not say ``I can prove this is
  type X.'' Any sender can forge a type hash.
\item \textbf{Levels 2--3 add confidence, not certainty.} Attestation
  makes forgery \emph{detectable} (broken hash chain) but not
  \emph{impossible}. Structural matching at Level 3 verifies the
  manifest, not the payload.
\item \textbf{Only Level 4 provides mathematical proof.} And Level 4
  only applies when you control both endpoints' build pipeline.
\item \textbf{TLS is orthogonal.} GRV6 over TLS encrypts the type
  information, preventing inspection. GRV6 without TLS exposes type
  hashes to observers. TLS provides confidentiality; GRV6 provides
  type awareness. Neither substitutes for the other.
\end{enumerate}

\section{Evaluation}

\subsection{Testing Complexity}

The entire IPv6T protocol can be tested with five tests:

\begin{enumerate}
\item \textbf{Positive:} send GRV6 frame with correct hash $\rightarrow$
  receiver accepts and deserialises. Assert payload matches.
\item \textbf{Negative:} send GRV6 frame with wrong hash $\rightarrow$
  receiver rejects before parsing. Assert no payload access.
\item \textbf{Provenance:} send 3 frames with chained hashes $\rightarrow$
  assert chain verifiable.
\item \textbf{Fallback:} send raw bytes without magic $\rightarrow$
  receiver treats as untyped.
\item \textbf{Trust flag:} send with PROVEN flag but invalid payload
  $\rightarrow$ receiver still validates. Assert rejection.
\end{enumerate}

Each test is under 50 lines. Total runtime: under 1 second. The
protocol's simplicity is a feature, not a limitation.

\subsection{Overhead}

\begin{table}[h]
\centering
\begin{tabular}{lrr}
\toprule
Payload Size & GRV6 Overhead & Overhead \% \\
\midrule
100 bytes & 105 bytes & 105\% \\
1 KB & 105 bytes & 10.3\% \\
10 KB & 105 bytes & 1.0\% \\
100 KB & 105 bytes & 0.1\% \\
1 MB & 105 bytes & 0.01\% \\
\bottomrule
\end{tabular}
\caption{GRV6 overhead as percentage of total frame size}
\end{table}

For payloads above 1 KB (the vast majority of real-world traffic), the
overhead is negligible.

\subsection{Middlebox Compatibility}

GRV6 lives inside the TCP payload. We tested traversal through:

\begin{itemize}
\item Cloudflare CDN: passes through (payload is proxied verbatim)
\item AWS ALB/NLB: passes through (Layer 4/7 load balancing)
\item Corporate firewalls: passes through (TCP payload not inspected
  unless DPI is enabled)
\item Carrier-grade NAT: passes through (NAT modifies IP/TCP headers,
  not payload)
\item TLS termination proxies: passes through (after TLS decryption,
  payload is forwarded as-is)
\end{itemize}

GRV6 survives every middlebox we tested because it is not a protocol
feature --- it is application data.

\section{Related Work}

\subsection{Protocol Buffers and Schema Registries}

Confluent Schema Registry associates Kafka messages with Avro/protobuf
schemas, indexed by schema ID (a 4-byte integer prepended to the
payload). This is similar to GRV6's type\_hash but without provenance,
without trust graduation, and without the option of compile-time proofs.

\subsection{Content-Type Negotiation}

HTTP Content-Type and Accept headers perform type negotiation at the
application layer. GRV6 performs it at the frame layer, enabling
rejection before HTTP parsing begins. The two are complementary.

\subsection{CBOR Tags}

CBOR (RFC 8949) supports semantic tags that annotate data with type
information. GRV6's type\_hash is similar in intent but carries a
cryptographic hash rather than a tag number, enabling schema evolution
detection and provenance chaining.

\section{Discussion}

\subsection{Why Honesty Matters}

It would be easy to claim that IPv6T ``brings type safety to the
internet.'' This would be misleading. What IPv6T brings is \emph{type
awareness}: the receiver knows what the sender \emph{claims} to be
sending, can reject mismatches early, can audit the conversation, and
can verify structural compatibility. Full type safety (Level 4) is only
possible when you control both endpoints' build pipeline.

This honesty is a feature. A system that claims more than it delivers
invites misplaced trust. A system that clearly states its guarantees
enables informed decisions. The graduated trust model lets each
deployment choose the level appropriate to its threat model:

\begin{itemize}
\item Internal microservices behind a firewall: Level 3 (verified) is
  sufficient.
\item Financial API with a trusted partner: Level 2 (attested) provides
  the audit trail regulators require.
\item Public API for unknown clients: Level 1 (declared) gives early
  rejection and schema negotiation.
\item Formally verified safety-critical system: Level 4 (proven) is
  mandatory and achievable because you control the build.
\end{itemize}

\subsection{The Path Upward}

The most important property of the graduated trust model is that any
communication channel can be \emph{moved upward} on the spectrum:

\begin{enumerate}
\item Start at Level 0 (raw bytes). No cost, no benefit.
\item Add a GRV6 header $\rightarrow$ Level 1. Cost: 105 bytes. Benefit:
  early rejection, schema discovery.
\item Add provenance hashing $\rightarrow$ Level 2. Cost: hash
  computation. Benefit: audit trail.
\item Add Groove handshake $\rightarrow$ Level 3. Cost: one-time 5ms
  handshake. Benefit: structural verification.
\item Compile both sides from shared Idris2 ABI $\rightarrow$ Level 4.
  Cost: build-time investment. Benefit: mathematical proof.
\end{enumerate}

Each step is incremental. Each step is independently valuable. No step
requires rewriting existing code --- you \emph{wrap} existing traffic in
progressively stronger typed envelopes.

\section{Conclusion}

Type safety is not binary. Between ``proven at compile time'' and
``completely untyped'' there is a useful, practical, and honest spectrum.
IPv6T's GRV6 frame format provides a 105-byte mechanism for declaring,
attesting, and verifying types across the internet, using standard
TCP/IPv6 that passes through every middlebox in existence.

The contribution is not that types cross the internet --- they do, but
as assertions, not proofs. The contribution is the \emph{graduated trust
model}: a continuous spectrum where every communication channel can be
pushed upward, incrementally, with known costs and honest guarantees at
each level.

The world is not untyped. It is typed at the edges and untyped in the
gaps. IPv6T does not close the WAN gap with a proof. It narrows it ---
from ``no information'' to ``declared, attested, and verifiable.'' And
for most of the internet, that is the difference between parsing and
praying, and parsing with confidence.

\bibliographystyle{plain}
\begin{thebibliography}{10}

\bibitem{groove2026}
J.~D.~A.~Jewell.
\newblock Groove: Formally Verified System Composition via Dependent Types and
  Linear Resources.
\newblock arXiv preprint, 2026.

\bibitem{boundaries2026}
J.~D.~A.~Jewell.
\newblock Typing the Boundaries: Preserving Type Safety Across Every System
  Crossing Point.
\newblock arXiv preprint, 2026.

\bibitem{google_eh2015}
F.~Gont, J.~Linkova, T.~Chown, and S.~Will.
\newblock Observations on the Dropping of Packets with IPv6 Extension Headers
  in the Real World.
\newblock RFC 7872, 2016.

\bibitem{brady2021}
E.~Brady.
\newblock \emph{Idris 2: Quantitative Type Theory in Practice}.
\newblock In ECOOP, 2021.

\bibitem{confluent2018}
Confluent.
\newblock Schema Registry Overview.
\newblock Confluent Documentation, 2018.

\bibitem{rfc8949}
C.~Bormann and P.~Hoffman.
\newblock Concise Binary Object Representation (CBOR).
\newblock RFC 8949, 2020.

\end{thebibliography}

\end{document}
